% Page 12 — Alban « la féodale »
% Fragment à intégrer au document principal avec \documentclass[12pt,a4paper]{article}

\usepackage{fontspec}
\usepackage[french]{babel}
\usepackage{geometry}
\usepackage{microtype}

\geometry{
  top=2.5cm,
  bottom=2.5cm,
  left=2.5cm,
  right=2.5cm
}

\setlength{\parindent}{1.2em}
\setlength{\parskip}{0.65em}

\begin{document}

\setcounter{page}{11}

\begin{center}
  {\Large\itshape Alban (au temps des Croisades)}
\end{center}

\vspace{1em}

L'époque des croisades a laissé -- grâce aux hospitaliers -- quelques souvenirs à Alban et dans son voisinage. Jusqu'à ces derniers temps -- dans l'enclos de St Léonard du Carmé (Montfranc) -- la tombe de l'Abbé Paul Routoulp, curé de Notre Dame d'Alban de 1763 à 1791, était surmontée d'une croix de pierre en forme de croix de Malte, insigne des Hospitaliers de St Jean de Jérusalem.

Mais un monument qui lui, est resté, c'est l'église de Notre Dame de la Garde. Une tradition orale dit qu'un pèlerin de Terre Sainte ait fait vœu de bâtir une chapelle à Notre Dame de la Garde, s'il revenait sain et sauf de son voyage après avoir essuyé maintes tempêtes. Revenu chez lui, il réalisait sa promesse. Et déjà en 1290, le pape Nicolas II accordait des indulgences à qui se rendait à ce sanctuaire. Les archives conservent le souvenir des pieuses démarches des Consuls d'Albi à la dévote chapelle à l'occasion de la peste de 1518. Au début du XIX\textsuperscript{ème} siècle, une grande dame, Louise de Bayne de Rayssac faisait à son amie Eugénie de Guérin, le récit de son pèlerinage matinal à Notre Dame de la Garde à la Pentecôte de 1837. Peut-être y eut-il de plus illustres pèlerins, mais nous l'ignorons.

\end{document}


\begin{center}
  {\Large\itshape Alban « la féodale »}
\end{center}

\vspace{1.5em}

\noindent
\lettrine{I}{l} est difficile d'établir les contours précis de l'Alban primitif.
En se référant à ce qui existait encore en 1598, nous pouvons faire une
reconstitution approximative d'une petite cité féodale entourée de murs
crénelés donnant sur des fossés extérieurs qui la protégeaient contre un
ennemi éventuel. Le centre en pouvait être l'actuelle place de l'église.
On accédait à notre bourg par la porte si l'on venait d'Albi, et par
« lou pourtal haut » si l'on arrivait de Montfranc. C'est non de la
« carreprolo » que se dressait dans une maison à trois étages la chambre
du guetteur. Il y avait encore en 1598 deux maisons à deux étages.
Les autres demeures devaient être modestes. Des groupes de jardins
« ourtalissos » jouxtaient les deux portails d'accès. On y pouvait voir
un « fourmial », un « grand carrairé », un « fieyrol ». La « fountanario »
mise à part, il est difficile de localiser la « fountetto si passado »,
loui mouli dal Bin, la « coumbo de bedel ». Les gens aujourd'hui ont
perdu le souvenir de bien des lieux aux noms charmants. A noter pour
les curieux et les poètes que fountaourio signifie « source d'or ».
Pourquoi ? La « carragrolo » est le chemin d'Alban à Peyreblanque,
près du pâtissier Rémy.
